% from ega0-1-fr.tex
\oldpage[0\textsubscript{I}]{76}

sur le système projectif d'anneaux $A/\mathfrak J$~; par passage à la limite
projective, on en déduit donc un module sur le séparé complété $\widehat A$
de $A$, que l'on appelle le \emph{produit tensoriel complété} de $M$ et de
$N$ et que l'on note $(M\otimes_A N)^\wedge$. Si l'on remarque que $M/V$
est canoniquement isomorphe à $\widehat M/\overline V$, où $\widehat M$ est
le séparé complété de $M$ et $\overline V$ l'adhérence dans $\widehat M$ de
l'image de $V$, on voit que le produit tensoriel complété
$(M\otimes_A N)^\wedge$ s'identifie canoniquement à
$(\widehat M\otimes_{\widehat A}\widehat N)^\wedge$, que l'on note aussi
$\widehat M\widehat\otimes_{\widehat A}\widehat N$.
\end{env}

\begin{env}[7.7.2]
\label{0.7.7.2-fr}
\begin{sloppypar}
Avec les notations précédentes, les produits tensoriels
$(M/V)\otimes_A(N/W)$ et
$(M/V)\otimes_{A/\mathfrak J}(N/W)$ s'identifient canoniquement~; ils
s'identifient donc aussi à
$(M\otimes_A N)/(\Im(V\otimes_A N)+\Im(M\otimes_A W))$.
On en conclut que $(M\otimes_A N)^\wedge$ est \emph{le séparé complété du
$A$-module $M\otimes_A N$, muni de la topologie pour laquelle les
sous-modules}
\[
  \Im(V\otimes_A N)+\Im(M\otimes_A W)
\]
\emph{forment un système fondamental de voisinages de $0$} ($V$ et $W$
parcourant l'ensemble des sous-modules ouverts de $M$ et $N$
respectivement)~; nous dirons pour abréger que cette topologie est le
\emph{produit tensoriel} des topologies données sur $M$ et $N$.
\end{sloppypar}
\end{env}

\begin{env}[7.7.3]
\label{0.7.7.3-fr}
Soient $M'$, $N'$ deux $A$-modules linéairement topologisés,
$u:M\to M'$, $v:N\to N'$ deux homomorphismes continus~; il est immédiat
que $u\otimes v$ est continu pour les topologies produits tensoriels sur
$M\otimes_A N$ et $M'\otimes_A N'$ respectivement~; par passage au séparé
complété, on en déduit un homomorphisme continu
$(M\otimes_A N)^\wedge\to(M'\otimes_A N')^\wedge$, que nous désignerons
par $u\widehat\otimes v$~; $(M\otimes_A N)^\wedge$ est donc un
\emph{bifoncteur} en $M$ et $N$ dans la catégorie des $A$-modules
linéairement topologisés.
\end{env}

\begin{env}[7.7.4]
\label{0.7.7.4-fr}
On définit de même le produit tensoriel complété d'un nombre fini quelconque
de $A$-modules linéairement topologisés~; il est immédiat que ce produit
possède les propriétés usuelles d'associativité et de commutativité.
\end{env}

\begin{env}[7.7.5]
\label{0.7.7.5-fr}
Si $B$, $C$ sont deux $A$-algèbres topologiques linéairement topologisées,
la topologie produit tensoriel sur $B\otimes_A C$ a pour système fondamental
de voisinages de $0$ les \emph{idéaux}
$\Im(\mathfrak K\otimes_A C)+\Im(B\otimes_A\mathfrak L)$ de l'algèbre
$B\otimes_A C$, $\mathfrak K$ (resp. $\mathfrak L$) parcourant l'ensemble
des idéaux ouverts de $B$ (resp. $C$). Par suite, $(B\otimes_A C)^\wedge$
est muni d'une structure de \emph{$\widehat A$-algèbre topologique}, limite
projective du système projectif de $(A/\mathfrak J)$-algèbres
$(B/\mathfrak K)\otimes_{A/\mathfrak J}(C/\mathfrak L)$ ($\mathfrak J$
idéal ouvert de $A$ tel que $\mathfrak J.B\subset\mathfrak K$,
$\mathfrak J.C\subset\mathfrak L$~; il en existe toujours). On dit que
cette algèbre est le \emph{produit tensoriel complété} des algèbres $B$ et
$C$.
\end{env}

\begin{env}[7.7.6]
\label{0.7.7.6-fr}
Les homomorphismes de $A$-algèbres $b\mapsto b\otimes 1$,
$c\mapsto 1\otimes c$ de $B$ et $C$ dans $B\otimes_A C$ sont continus
quand on munit cette dernière algèbre de la topologie produit tensoriel~;
par composition avec l'homomorphisme canonique de $B\otimes_A C$ dans son
séparé complété, ils donnent donc des \emph{homomorphismes canoniques}
$\rho:B\to(B\otimes_A C)^\wedge$, $\sigma:C\to(B\otimes_A C)^\wedge$.
L'algèbre $(B\otimes_A C)^\wedge$ et les homomorphismes $\rho$ et $\sigma$
possèdent en outre la \emph{propriété universelle suivante}~: pour toute
$A$-algèbre linéairement topologisée, séparée et complète $D$, et tout
couple de $A$-homomorphismes continus $u:B\to D$, $v:C\to D$, il existe un
$A$-homomorphisme continu et un seul
$w:(B\otimes_A C)^\wedge\to D$ et un seul tel que $u=w\circ\rho$ et
$v=w\circ\sigma$. En effet, il existe déjà un $A$-homomorphisme unique
$w_0:B\otimes_A C\to D$ tel que $u(b)=w_0(b\otimes 1)$ et
$v(c)=w_0(1\otimes c)$, et tout revient à prouver que $w_0$ est continu,
car il donnera alors un

